\section*{Resumo}
Este relatório descreve o projeto de uma unidade de escrita e leitura
multiprotocolo que liga um PC a dispositivos externos através de I2C,
RS232, RS485, dois barramentos CAN FD e Ethernet, e que comanda três
saídas PWM isoladas com marcação temporal sincronizada por PPS/GPS. A
unidade encaminha dados entre interfaces com taxas diferentes, mantém
progresso sob concorrência e guarda a configuração em memória não
volátil.

A partir das necessidades recolhidas junto dos docentes, definem-se
sete casos de utilização, 18 requisitos funcionais, 8 requisitos não
funcionais e os respetivos critérios de aceitação, e propõe-se uma
arquitetura de alto nível. Na fase de projeto escolhe-se o
microcontrolador STM32H743 com FreeRTOS, transceivers isolados para CAN FD,
RS485 e PWM, e uma fonte buck-boost que aceita de 5 a 60 V. O esquema
elétrico completo é gerado por código e verificado automaticamente, e o
firmware organiza-se em tarefas com filas por interface e duas classes de
prioridade.

\section*{Palavras-chave}
Gateway multiprotocolo, CAN FD, RS485, Ethernet, PPS/GPS, PWM isolado,
STM32H743, FreeRTOS.
